% Auto-extracted body from wbc_section_draft.tex for standalone preview
% Cross-refs to main.tex labels replaced with text for standalone compilation

\subsection{Whole-Body Control Baselines}
\label{sec:wbc}

To address whether a tuned whole-body controller (WBC) could match ACC, we evaluated two WBC configurations that were developed and subsequently discarded during controller development:
\textbf{(W1)}~a task-priority quadratic-program (QP) WBC---minimizing COM height error, torso orientation error, posture damping, and wheel regularization subject to full-body dynamics and contact constraints---used as the \emph{primary} balance controller; and
\textbf{(W2)}~the same WBC blended as an \emph{assist} on top of ACC via $\boldsymbol{\tau}_{\text{assist}} = \boldsymbol{\tau}_{\text{ACC}} + \alpha\!\cdot\!(\boldsymbol{\tau}_{\text{wbc}} - \boldsymbol{\tau}_{\text{ACC}})$, with a seven-term multiplicative gate $\alpha_{\!j} = \alpha_{\max}\!\cdot\!g_{\text{stability}}\!\cdot\!g_{\text{height}}\!\cdot\!g_{\text{push}}\!\cdot\!g_{\text{divergence}}\!\cdot\!A_j\!\cdot\!K_{\text{role},j}$ that continuously modulates per-joint WBC authority.
Both were evaluated under identical cloned simulation conditions (three-arm counterfactual: ACC baseline, WBC-only, ACC$+$WBC assist; 225 scenarios spanning fixed-height balance, height transitions, and push recovery at \SIrange{20}{120}{N}; ACC profile \texttt{K2\_JAX\_DEDICATED\_DEFAULT\_V3}; no gain retuning).

\medskip\noindent\textbf{W1 (WBC as primary controller).}
The standalone QP-WBC fails catastrophically: \textbf{\num{17.2}$\times$} more falls than ACC (5{,}985 vs.\ 347 across 225 scenarios), \textbf{\num{28.4}$\times$} more single-push falls (3{,}353 vs.\ 118), and pitch RMS \SI{41.8}{\degree} vs.\ ACC's \SI{15.9}{\degree} ($\num{2.63}\times$).
Survival time at \SI{0.60}{m} is \SI{0.52}{s} (\SI{100}{\percent} fall rate).
The failure mechanism is structural: the WBC's posture task is a velocity damper ($\ddot{\mathbf{q}}_{\text{joints}} = -k_{\text{posture}}\!\cdot\!(\mathbf{q} - \mathbf{q}_{\text{current}})$, $k_{\text{posture}}{=}10$), not a position tracker---it cannot hold a target posture.
It lacks gravity-compensation feedforward, uses a single linearization at \SI{0.67}{m} (vs.\ the \SIrange{0.45}{0.60}{m} operating range), and treats wheels as point contacts without a rolling constraint.
The QP solver fails on \SI{14.2}{\percent} of steps, creating torque discontinuities.
Four critical implementation bugs were identified post-hoc (contact Jacobian zero-gradient, solver feasibility-only hardcode, base-z vs.\ CoM-z height reference mismatch, and gate height-reference conflict; detailed in Supplementary).
WBC as-implemented-and-tuned therefore fails as a standalone controller.
Critically, its failure is not solely a tuning or implementation defect: the WBC lacks the \emph{regime-gating} that ACC's proximity-gated anchor provides.
A WBC distributes torque optimally for a fixed task set---it cannot express qualitatively distinct physical regimes (quiet stance, push ringdown, large excursion) as distinct control modes.
The precision--recovery trade-off ACC resolves through spatial gating and asymmetric temporal envelope-following is not expressible as a task-priority QP objective; even a bug-free WBC would lack the structured prior that gates the anchor integral, and would therefore face the same global-integral windup that our ablation S4 (no proximity gate, see push ablation table in the main text) confirms is fatal for sub-millimeter precision.

\medskip\noindent\textbf{W2 (WBC as assist on ACC).}
The assist gate---designed to be maximally conservative---blocks WBC contribution in ${\ge}\SI{99}{\percent}$ of steps:
in the 225-scenario batch, all metric ratios are \textbf{\num{1.000000}} (ACC$+$WBC assist $\equiv$ ACC within measurement noise); a 48-scenario quick evaluation classifies all 48 scenarios as \texttt{ASSIST\_EQUIVALENT} (zero regressions, zero safety violations).
This zero-net effect is \emph{independent} of the WBC implementation bugs: the gate blocks WBC \emph{before} its torque is computed, based on physical state (pitch, roll, height deviation, push magnitude, divergence), so the WBC's internal correctness is irrelevant to the assist outcome.

When gates are manually overridden to force WBC contribution (keyframe test at \SI{0.53}{m}, $\alpha_{\max}{=}0.50$, widened stability and height thresholds), the assist produces mixed results:
pitch RMS improves \SI{17}{\percent} (\SI{2.44}{\degree} $\rightarrow$ \SI{2.04}{\degree}), drift improves \SI{18}{\percent} (\SI{0.070}{m} $\rightarrow$ \SI{0.058}{m}), and yaw drift improves \SI{20}{\percent}---but roll RMS \emph{degrades} \SI{35}{\percent} (\SI{0.20}{\degree} $\rightarrow$ \SI{0.27}{\degree}).
The degradation arises from a fundamental architectural conflict:
ACC's lateral roll controller and WBC's torso-orientation task both write to hip-roll joints, but from incompatible philosophies (decentralized PD vs.\ centralized QP exploiting full mass-matrix coupling).
A posture-guided variant (WBC outputs posture targets $\mathbf{q}_{\text{ref}}$, ACC executes torque) reduced the roll degradation from $+$\SI{35}{\percent} to $+$\SI{7}{\percent} by eliminating torque-level interference, but did not produce a net improvement over ACC alone.

\medskip\noindent\textbf{Timing overhead.}
The WBC QP solve alone consumes \SI{0.45}{ms} mean, \SI{2.94}{ms} p95, and \textbf{\SI{3.43}{ms} p99} (OSQP, offline Python/NumPy path, 4 scenarios).
The p99 solve time alone exceeds the \SI{0.5}{ms} end-to-end budget margin below the \SI{10}{ms} performance threshold by $\num{6.9}\times$---before accounting for matrix assembly, contact Jacobian computation, or the ACC computation baseline.
An incremental JAX-JIT path was prototyped to reduce per-step overhead but was never integrated into the real-time control loop; real-time WBC assist overhead therefore remains uncharacterized.
Even if the solve time could be reduced below the margin, the assist adds no measurable balance benefit (W2 zero-net result), so the compute cost is pure overhead.

\medskip\noindent\textbf{Summary.}
Neither standalone WBC nor WBC-as-assist resolves the precision--recovery trade-off on this platform.
The trade-off is resolved by regime-gating---the proximity-gated anchor and asymmetric envelope follower---which WBC's task-priority formulation does not express.
ACC's conditional gates therefore provide the essential structured prior; the WBC gap is one of structural expressiveness, not torque optimality.

% --- Supplementary section ---
\subsection{Supplementary: WBC Implementation Details}
\label{sec:wbc_supp}

\subsubsection{W1 Standalone WBC — Full Metric Table}

Table~\ref{tab:wbc_w1_full} reports all measured metrics for the standalone QP-WBC (W1) vs.\ ACC, from the 225-scenario three-arm counterfactual evaluation and the 4-scenario Phase~3D full-batch execution.
The WBC was evaluated with known implementation bugs (Table~\ref{tab:wbc_bugs}); W1 numbers may therefore be pessimistic.
However, the failure mechanism---velocity-damping posture, single-linearization, no rolling model---is architectural, and the absence of regime-gating (the core ACC contribution) is structural, not bug-dependent.

\begin{table}[h]
\centering
\caption{Standalone QP-WBC (W1) vs.\ ACC — Full Metric Comparison}
\label{tab:wbc_w1_full}
\footnotesize
\begin{tabular}{@{}lccc@{}}
\toprule
Metric & WBC-only & ACC (V3) & Ratio (WBC/ACC) \\
\midrule
Total falls (225 scenarios)                                  & 5{,}985   & 347   & \textbf{17.2}$\times$ \\
Single-push falls                                            & 3{,}353   & 118   & \textbf{28.4}$\times$ \\
Falls (Phase~3D, 4 step\_e scenarios)                        & 14{,}713  & 0     & --- \\
Survival time at \SI{0.60}{m} (s)                            & 0.52      & 20.0  & 0.026$\times$ \\
Pitch RMS (\si{\degree})                                     & 41.8      & 15.9  & 2.63$\times$ \\
Roll RMS (\si{\degree})                                      & 96.8      & 84.2  & 1.15$\times$ \\
Yaw drift (\si{\degree})                                     & 8.0       & 4.24  & 1.89$\times$ \\
Height RMS (m)                                               & 0.253     & 0.261 & 0.97$\times$ \\
Planar drift (m)                                             & 2.57      & 2.42  & 1.06$\times$ \\
Posture error (ratio vs.\ ACC)                               & ---       & ---   & \textbf{8.05}$\times$ \\
Yaw error (ratio vs.\ ACC)                                   & ---       & ---   & \textbf{15.43}$\times$ \\
QP solver success rate                                       & \SI{85.85}{\percent} & --- & --- \\
\bottomrule
\end{tabular}
\vspace{4pt}\par\noindent\footnotesize
Sources: 225-scenario batch from \texttt{V3\_vs\_V3\_Assist\_comparison\_report.md};
Phase~3D from \texttt{full\_batch\_verdict.json};
survival time from \texttt{best\_torque\_wbc\_summary.json}.
All runs use ACC profile, no gain retuning.
WBC evaluated with known bugs (Table~\ref{tab:wbc_bugs}); W1 numbers may be pessimistic.
\end{table}

\subsubsection{Known WBC Implementation Bugs (F1--F4)}

Table~\ref{tab:wbc_bugs} lists the four critical bugs identified during post-hoc audit of the WBC implementation.
These bugs affect W1 (standalone) metrics; W2 (assist) zero-net is bug-independent because the gate blocks WBC torque before it is computed.

\begin{table}[h]
\centering
\caption{WBC Implementation Bugs (Post-Hoc Audit)}
\label{tab:wbc_bugs}
\footnotesize
\begin{tabular}{@{}c>{\raggedright\arraybackslash}p{0.35\linewidth}>{\raggedright\arraybackslash}p{0.40\linewidth}@{}}
\toprule
ID & Description & Effect on Metrics \\
\midrule
F1 & Contact Jacobian zero-gradient: all leg joint columns~$=0$ (autodiff through cached FK). &
Ground reaction forces produce zero moment on leg joints in QP dynamics; $\boldsymbol{\tau}_{\text{wbc}}$ meaningless at hip-pitch/knee ($K_{\text{role}}{=}0.60$). \\
F2 & Fast solver (OSQP, default) hardcodes \texttt{feasibility\_only}---drops entire task stack. &
QP solves for minimum-norm $\ddot{\mathbf{q}}$ with zero task objectives; all task modes produce identical output. \\
F3 & Offline eval passes base-$z$ (\SI{0.536}{m}) as commanded height instead of CoM-$z$ (\SI{0.404}{m})---\SI{13.2}{cm} error. &
ACC height-dependent gain schedule, feedforward, and equilibrium references at wrong operating point; drift/heading/centering gates disabled. \\
F4 & Gate adaptive assist height reference conflict: WBC model at \SI{0.67}{m}, V3 commands CoM at \SI{0.45}{--}\SI{0.60}{m}. &
$g_{\text{height}} \approx 0$ at all test heights; WBC gate forced closed by multiplicative safety floor. \\
\bottomrule
\end{tabular}
\vspace{4pt}\par\noindent\footnotesize
Source: \texttt{docs/validation/v3\_wbc\_assist\_root\_cause\_audit.md}.
F1--F4 affect W1 (standalone); architectural argument (absence of regime-gating) does not depend on them.
W2 (assist) zero-net is bug-independent: gate blocks WBC before torque is computed.
\end{table}

\subsubsection{W2 Zero-Net is Bug-Independent}
The assist gate $\alpha_{\!j}$ is computed from physical state (pitch, roll, height deviation, push magnitude, divergence) \emph{before} WBC torque is evaluated.
Since $\alpha_{\!j} \approx 0$ in ${\ge}\SI{99}{\percent}$ of steps, the WBC's internal correctness---buggy or not---is irrelevant to the assist outcome.
The 225-scenario equivalence ($\text{ACC}+\text{WBC assist} \equiv \text{ACC}$ within measurement noise) is therefore a robust result.
The keyframe test (gates manually overridden) reveals the torque-blend conflict (roll~$+$\SI{35}{\percent}) independently of F1--F4: with the gate forced open, WBC torque is applied, and the decentralized-vs-centralized conflict manifests as roll degradation.
